\documentclass[a4paper, 10pt]{article}

\usepackage[utf8]{inputenc}
\include{page-style}

%%% Title
\renewcommand{\maketitle}{
    \begin{center}
        {\huge\scshape Loan~Repayment~Calculations}
        \rule{\textwidth}{1.6pt}\vspace*{-\baselineskip}\vspace*{2pt}
        \rule{\textwidth}{0.4pt}\\[0.2\baselineskip]
        {\Large\quad\quad William~Wallis\hfill\today \quad\quad}
    \end{center}
}



\begin{document}
\pagestyle{plain}

\maketitle
\tableofcontents

\section{Definitions}\label{sec:definitions}
\subsection{Notation}\label{subsec:notation}

A loan is a fixed value of money borrowed by an entity and usually repaid over a series of instalments.\ The following notation is used for each of the components of the loan:
\begin{itemize}
    \item[$L$:] Loan amount.
    \item[$R$:] Periodic interest rate.
    \item[$N$:] Total number of repayments.
    \item[$P$:] Total periodic repayment value.\ If this changes over time, the value at period $n$ is denoted $P_{n}$.
    \item[$b$:] Whether the interest is applied before or after the repayment.
    \item[$B_{n}$:] The balance on the loan at period $n$.
\end{itemize}

The periodic repayment for a loan usually has (at least) 2 components:
\begin{itemize}
    \item[$P_{P}$:] The principal part of the periodic repayment value, which is paying off the original money that was borrowed.
    \item[$P_{I}$:] The interest part of the periodic repayment value, which is paying off the interest applied on the loan.
\end{itemize}

In `real life', a loan can have other components such as fees.\ These are outside the scope of these calculations.


\subsection{Loan~Types}\label{subsec:loan-types}


\section{Periodic~Repayment~Amount}\label{sec:periodic-repayment-amount}


\end{document}
